\section{Fifth axis: the $(\Eavail,W)$ deep-inelastic-tail study}
\label{sec:eavailw}

Enabling Valencia 2p2h resolves the low-recoil excess (\S\ref{sec:3d-2p2h})
but leaves a separate high-$\Eavail$ data excess where 2p2h does not
contribute and the sub-percent pion-FSI dials cannot reach. A corrected
$(\Eavail,W)$ covariance is now adopted --- under a digest-bound exception and
carrying four measurements that qualify what it supports
(\S\ref{sec:eavailw-adopted}) --- so the comparison below, which was made at
central-value level while none existed, is stated together with that
covariance and the conditions on it. GiBUU spreads its deficit across the entire
tail (data/generator $1.59$ at $\Eavail\in[0.8,1.5)$, $1.36$ at $[1.5,3.0)$)
rather than piling it in the deep catch bin like the GENIE variants --- an
independent, qualitatively distinct confirmation. To test the deep-inelastic
hypothesis directly, we added the hadronic invariant mass $W$ as a fifth
OmniFold axis.

\subsection{$W$ as a fifth axis}
\label{sec:eavailw-axis}
Truth $W$ is the experimenter's invariant mass
$W=\sqrt{M^{2}+2M(E_\nu-E_\mu)-Q^{2}}$ with $M$ the struck-nucleon mass
(\texttt{CVUniverse::GetTrueExperimentersW()}); reco $W$ uses the same
calorimetric energy transfer $q_{0}$ and muon-kinematic $Q^{2}$ as the reco
$q_{3}$ estimator, with the average nucleon mass
(\texttt{CVUniverse::RecoW()}). Like $q_{3}$ --- and unlike $\Eavail$ --- $W$
is not invariant under the lateral systematic shifts, so the event loop dumps
shifted-$W$ lateral universes per universe, mirroring the $q_{3}$ treatment
(\S\ref{sec:4d}); the dump also carries truth hadronic diagnostics (proton and
pion multiplicities, leading-hadron angle). The $W$
binning $[0,1.1,1.4,1.8,2.2,3.0,100]\,\si{GeV}$ brackets the quasielastic
peak, the $\Delta$, the transition region, and the DIS tail. The five-axis
unfold $(\pT,\pPar,\Eavail,q_{3},W)$ passes the same marginalization anchors
as its predecessors: the $W$-marginal recovers the frozen 4D cross section to
\SI{\eavailMargin}{\percent} in normalization (5D/4D total $=1.0011$), with per-axis
median differences of \SIrange{0.3}{1.5}{\percent}, and the
$d\sigma/dW$ spectrum is finite and non-negative across all six bins.

\subsection{Where the excess lives in the $(\Eavail,W)$ plane}
\label{sec:eavailw-map}
Projecting the unfolded five-axis result onto $(\Eavail,W)$ and subtracting
the MINERvA Tune~v1 prediction (the analysis simulation's truth weighted by its full tune weight \texttt{w\_truth}, i.e.\ the OmniFold prior; \emph{not} the untuned GENIE~CV of \S\ref{sec:eavailw-band}) localizes the excess
(Fig.~\ref{fig:eavailWexcess}): it is \emph{predominantly deep-inelastic}.
The two highest-$\Eavail$ bands carry \SI{57}{\percent} of the positive
excess; of all positive excess, high-$\Eavail$ ($\ge\SI{0.8}{GeV}$) holds
\SI{67}{\percent}, and \SI{83}{\percent} of \emph{that} sits at
$W\ge\SI{1.8}{GeV}$, with the single largest cell the deep-DIS corner
($\Eavail>\SI{3}{GeV}$, $W>\SI{3}{GeV}$) at \SI{22}{\percent}. A secondary
low-$W$ (below \SI{1.1}{GeV}, QE-like) excess and a small $\Delta$-region
($W\simeq1.4$--$\SI{1.8}{GeV}$) deficit are also present.
This localization is a central-value statement under the nominal detector
response. The shares above are central-value quantities read off the projected
five-axis result, and that result is not seed-independent: on the 43 projection functionals of the estimator-seed campaign (these
\zprojCells\ destinations plus the total rate), the estimator seed moves the central value by at most
\SI{\cvMoveProjMax}{\percent}, which is at most
\SI{\cvMoveOverSigmaMax}{\percent} of the quoted uncertainty (median
\SI{\cvMoveOverSigmaMed}{\percent}); see \S\ref{sec:eavailw-adopted}. That
bound is on the projected central values themselves. The shares are ratios of
differences formed from them and are not separately measured against the seed.
The response-mismatch closure in
Appendix~\ref{app:response-mismatch} is designed to test whether an incorrect
hadronic response can move strength into or out of the low-$W$ component; it
has not yet produced a result.

\begin{figure}[H]
  \centering
  \includegraphics[width=0.85\textwidth]{excess_eavail_W}
  \caption{Unfolded data versus the MINERvA Tune~v1 prediction (GENIE~2.12.6 with the analysis tune weights \texttt{w\_truth}, the OmniFold prior; labelled ``GENIE CV'' and ``CV'' inside the figure) across the
  $(\Eavail,W)$ plane. \emph{Left:} $d\sigma/d\Eavail$ for both, integrated over
  $W$. \emph{Center:} the per-cell ratio data/CV, dimensionless, on a
  $0.5$--$1.5$ colour scale. \emph{Right:} the per-cell difference
  data\,$-$\,CV.
  \textbf{Axes are cell indices, counting from zero, on the
  $7\times6$ reporting grid}: $\Eavail$ edges
  $[0,0.1,0.2,0.4,0.8,1.5,3.0,100]\,\si{GeV}$ and $W$ edges
  $[0,1.1,1.4,1.8,2.2,3.0,100]\,\si{GeV}$, so index~6 on $\Eavail$ and
  index~5 on $W$ are wide catch bins.
  \textbf{Units and the bin-area factor:} the centre panel is a ratio of
  densities and is area-free, but the right panel is the difference of
  \emph{cell-integrated} cross sections,
  $(\mathrm{d}^{2}\sigma/\mathrm{d}\Eavail\,\mathrm{d}W)\times
  \Delta\Eavail\,\Delta W$ in \si{cm^2/nucleon}
  (\texttt{nd-unfolding/excess\_eavail\_W.py}), so the catch-bin cells carry
  their large widths and look correspondingly bigger. Read the two panels
  together: the positive excess concentrates at high $\Eavail$ \emph{and} high
  $W$ --- the deep-inelastic region --- with a secondary low-$W$ (QE-like)
  component and a small $\Delta$-region deficit. The percentage shares quoted in
  the text are shares of the cell-integrated positive excess, i.e.\ of the right
  panel.}
  \label{fig:eavailWexcess}
\end{figure}

\subsection{Generator band: the corner is generator- and tune-robust}
\label{sec:eavailw-band}
We confirm this against independent spline-normalized $(\Eavail,W)$
predictions from GENIE~CV, GENIE$+$Valencia~2p2h/MEC, NuWro, and GiBUU
(Fig.~\ref{fig:eavailWband}). These are separate productions from the 3D samples of \S\ref{sec:3d-models}; this GENIE~CV's in-phase-space total, \SI{2.44e-38}{cm^2/nucleon}, is \SI{3}{\percent} below the \SI{2.52e-38}{} quoted there, and that difference has not been reconciled.
% --- PROVENANCE / REDUCTION for the corner ratios (J27.3, recorded 2026-08-02) ---
% The ratios below are CORNER-INTEGRATED cross-section ratios -- not per-bin means,
% not a weighted mean, not the worst bin. Reduction, from
% 3d-unfolding/genie/overlay_eavailW_band.py:88-108:
%   corner = outer(EAVAIL_EDGES[:-1] >= 0.8, W_EDGES[:-1] >= 1.8)  -> 3x3 = 9 cells
%            [Eavail 0.8-1.5, 1.5-3.0, 3.0-100] x [W 1.8-2.2, 2.2-3.0, 3.0-100]
%   data/gen = sum_corner(hData2D * dEavail * dW)
%              / sum_corner(hXSec_eavailW * dEavail * dW)
% Source of 1.54/1.58/1.56: ND_OMNIFOLD_RUN_LOG.md:988-990 (2026-06-08 band run),
% which reads "All THREE underpredict ... 1.54 CV, 1.58 +MEC, 1.56 NuWro" and states
% "GiBUU excluded (FinalEvents.dat lacks per-event Enu)". GiBUU was structurally
% ABSENT from the run these three numbers come from. This paragraph previously read
% "All four underpredict ... by 54-58%" -- a later edit that attached three
% magnitudes to four generators without a recomputation; corrected 2026-08-02 to
% attribute them to the three that produced them. GiBUU was afterwards added to
% the band (gibuu_to_xsec_eavailW.py, Task 16; make_figures.sh:55 passes it and the
% figure legend shows it), so the 4-generator overlay DID print a GiBUU corner
% ratio -- to stdout, never captured. To close: re-run make_figures.sh:55 and read
% the "hiE-hiW corner ... data/gen=" line. Inputs live only on Perlmutter /pscratch.
% Do NOT substitute either of these -- both are different reductions:
%   - Eavail-projected GiBUU data/gen 1.59 / 1.36 / 1.91 (RUN_LOG:1126-1127)
%     integrates over ALL W, not W>=1.8.
%   - corner chi2 381.1/12 (RUN_LOG:1155-1157) uses Eavail>=0.4 -> 12 cells,
%     a different sub-block that adds the 0.4-0.8 Eavail row.
% Until recomputed, 54-58% is a THREE-generator span. When the GiBUU ratio arrives
% this is a one-line upgrade: add it, and widen the span only if it lands outside.
%
% CLOSED 2026-08-11. All FOUR ratios recomputed TOGETHER on one data file, so the set
% is internally consistent rather than a fourth number appended to three older ones.
% Log (whole stream, no tail): 3d-unfolding/genie/eavailW_band_20260811_allfour.log
% Ledger: "2026-08-11 (E_avail,W) four-generator corner ratios". Predeclared at
% docs/orchestration/PREDECLARE-20260811-gibuu-corner-ratio.md, branch G1.
%   GENIE-CV  corner 8.7918e-39  data/gen 1.535    (note had 1.54)
%   GENIE+MEC corner 8.5484e-39  data/gen 1.579    (note had 1.58)
%   NuWro     corner 8.6369e-39  data/gen 1.563    (note had 1.56)
%   GiBUU     corner 8.3893e-39  data/gen 1.609    <- NEW; lands OUTSIDE 1.54-1.58,
%   data      corner 1.3497e-38                       so the span widens 54-58 -> 54-61
% The three reproduce AT THE PRECISION THIS NOTE QUOTES. Identity is NOT established:
% the only surviving record of the 2026-06-08 values is RUN_LOG:988-990 at three
% significant figures, and the data file (excess_eavail_W.root, 2026-07-14) POSTDATES
% them by five weeks -- so a shift in the third decimal is masked by rounding. Per
% BEN-086, agreement at printed precision is not identity; do not re-quote these at
% more digits from the 2026-06-08 lineage.
% Controls that reproduce EXACTLY and are why the set is trustworthy: GiBUU integrated
% 2.2227e-38 and data 3.0699e-38 (this paragraph's own 2.22 / 3.07). The ORDERING
% GiBUU < NuWro < GENIE CV matches sec_3d.tex:215-217, but the GENIE value does not:
% this band's GENIE-CV 2.4446e-38 (work_seed* gevgen, gen_to_xsec_eavailW.py) is a
% different object from sec_3d's 2.52e-38 (Stage-A work_p* gevgen, genie_to_xsec3d.py);
% the 2.95% difference is not reconciled. (The ledger's VL35 prose repeats the old claim.)
%
% NORMALIZATION INCONSISTENCY, FOUND WHILE CLOSING THIS AND DELIBERATELY LEFT ALONE.
% This paragraph states two deficits, two sentences apart, in TWO DIFFERENT
% NORMALIZATIONS, both rendered as a bare percentage with no marker:
%   "underpredict ... by 54-61%"  is generator-relative, (data-gen)/gen = data/gen - 1
%   "sit 23-26% below the data"   is data-relative,      (data-gen)/data
% Both are internally correct and each matches its own arithmetic. But in a COMMON
% normalization the corner deficit is 34.9-37.8% data-relative (or the high-W deficit
% is 29.8-34.8% generator-relative), so a reader comparing "54-61%" against "23-26%"
% infers a contrast of ~2.4x where the consistent answer is ~1.2x. Choosing the
% convention is an authorial decision about a physics claim, so it was ROUTED rather
% than silently changed; the ratios above are unaffected either way.
%
% RESOLVED 2026-08-12 by Joseph (-> Session A -> Session B, BEN-082(v)): "use one
% normalization throughout -- prefer explicit data/generator ratios for both
% comparisons rather than mixing 'percent above generator' with 'percent below data.'"
% Both sentences now state data/generator and neither states a bare percentage.
% The W-band ratios were COMPUTED FROM THE SOURCE CROSS-SECTIONS, not converted from
% the rounded 23-26%, because converting rounded endpoints would manufacture precision
% the source does not have (the BEN-086 trap this same comment block warns about).
% From eavailW_band_20260811_allfour.log:22-23, W in [2.2,3.0) GeV, units cm^2/nucleon:
%   data 7.481e-39
%   GENIE-CV  5.764e-39 -> 1.2979    GENIE+MEC 5.622e-39 -> 1.3307
%   NuWro     5.715e-39 -> 1.3090    GiBUU     5.550e-39 -> 1.3479
%   span 1.298-1.348 -> quoted 1.30-1.35
% CROSS-CHECK, and it is why these are trustworthy: the same operands reproduce the
% note's previous data-relative wording exactly -- 22.95%, 24.85%, 23.61%, 25.81%,
% i.e. "23-26% below the data". The prose changed convention; no number moved.
% The corrected contrast is now legible on the page: 1.54-1.61 against 1.30-1.35.
% --------------------------------------------------------------------------------
GENIE~CV, GENIE$+$MEC, NuWro, and GiBUU all underpredict the
high-$\Eavail\times$high-$W$ corner, with data/generator $=1.54$, $1.58$,
$1.56$, and $1.61$ respectively (span $1.54$--$1.61$). Crucially,
enabling Valencia 2p2h does \emph{not} close the corner --- it slightly
worsens it ($1.54\to1.58$), since 2p2h populates low $W$ --- and the
independent NuWro and GiBUU generators miss it by the same margin or more.
These are central-value ratios. A projected covariance now exists
(\S\ref{sec:eavailw-adopted}); the compatibility metric it would support is
\emph{not} quoted here, because the estimator-seed sensitivity measured on that
covariance exceeds the movement limit declared for it before production, and a
$\chi^{2}$ computed against a matrix carrying an unclosed sensitivity of that
size would state a precision the input does not have.
GiBUU is the most deficient generator both overall (integrated
$\sigma=\SI{2.22e-38}{cm^2/nucleon}$ against the data's \SI{3.07e-38}{}) and
in the corner, where it sets the upper end of the band; no generator
considered here reaches the corner. At
$W\in[2.2,3.0)\,\si{GeV}$ the same four generators underpredict by a smaller
factor, data/generator $=1.30$, $1.33$, $1.31$, and $1.35$ (span
$1.30$--$1.35$).

\begin{figure}[H]
  \centering
  \includegraphics[width=0.99\textwidth]{eavailW_band}
  \caption{Unfolded $(\Eavail,W)$ cross section versus the four-generator
  band (GENIE~CV, GENIE$+$MEC, NuWro, GiBUU). Every generator underpredicts
  the high-$W$, high-$\Eavail$ DIS corner; adding 2p2h does not help there
  because 2p2h populates low $W$.}
  \label{fig:eavailWband}
\end{figure}

Dedicated MINERvA measurements of shallow-inelastic neutrino and antineutrino
scattering also report shape and normalization disagreements with generators
\cite{MINERvA:2025sis}. They motivate tests across the resonance--DIS transition,
but their selections and reported observables differ from this inclusive
projection. They therefore provide physics context, not an independent
confirmation of the same localized excess. The location identified in the
central-value map is descriptive; a future significance must account for any
selection of the tested region using these data.

\subsection{The earlier direct $(\Eavail,W)$ block sum, and what its
replacement had to satisfy}
\label{sec:eavailw-cov}
\noindent\textbf{This subsection is historical. It describes a superseded
construction and the requirements imposed on its replacement; the covariance
this note actually quotes is specified in \S\ref{sec:eavailw-construction}.}

An earlier attempt built a dedicated 42-bin $(\Eavail,W)$ covariance
\emph{directly}, as $C = C_{\mathrm{syst}}+C_{\mathrm{stat}}+C_{\mathrm{lateral}}$,
by a \emph{frozen-reweighter} block sum on the merged five-axis
\texttt{\_universes\_full} sample: the central-value OmniFold push-weights
were held fixed and only the per-event truth reweight ratios changed per
systematic universe (the same frozen-reweighter \emph{technique} as the PET
budget, \S\ref{sec:pet-syst} --- not the PET covariance matrix, which is never
applied to this GBDT central value; no per-universe re-inference). That object
carried 13 knob bands $\pm1\sigma$, the 100-universe flux multisim, the
projected full five-axis bootstrap covariance $C_{EW}=M C_{5D}M^{\top}$, and a
$W$-resolved detector (lateral) block. It is not adopted and is not quoted
anywhere in this note.

Three requirements were declared for any replacement, and all three are met by
the adopted object of \S\ref{sec:eavailw-construction}: project the corrected
five-axis statistical covariance as $M C_{5D}M^{\top}$ rather than rebuilding it
in the destination binning; use the corrected systematic endpoint convention;
and bound the migrations omitted by the CV-support-limited dump-all sample with
selection-complete active-universe event loops. That last replacement has
landed. Its detector block uses five shifted-kinematic bands
(\texttt{Muon\_Energy\_MINERvA/MINOS}, \texttt{MuonResolution}, and
\texttt{BeamAngleX/Y}) plus weight-only \texttt{MinosEfficiency} and GEANT
bands; the scope that classification does and does not establish is stated in
\S\ref{sec:systematics}.

\subsection{The adopted construction, in full}
\label{sec:eavailw-construction}
\noindent\textbf{This subsection is the authoritative account of the covariance
this note quotes.} Where an earlier recipe appears elsewhere in this note
(\S\ref{sec:eavailw-cov} for the direct $(\Eavail,W)$ block sum,
\S\ref{app:pair} for the per-band pair convention), that recipe is historical
or partial and this one governs.

\paragraph{The five-axis object.} On the $\latBins$ reported cells of the
$(\pT,\pPar,\Eavail,q_{3},W)$ grid, the adopted covariance is
\begin{equation}
  \label{eq:zassembly}
  C_{Z} \;=\; D\,\Big(\textstyle\sum_{V} C_{b}\Big)\,D
  \;+\; \textstyle\sum_{R} C_{b}
  \;+\; \textstyle\sum_{A} L_{b}
  \;+\; C_{\mathrm{stat}}
  \;+\; C_{\mathrm{ML}},
\end{equation}
with the inflation $D=\diag(g)$,
\begin{equation}
  \label{eq:zinflation}
  g_{i} \;=\;
  \sqrt{\frac{\max\!\big(v^{\mathrm{cv}}_{\mathrm{uni},i},\,
                         v_{\mathrm{blk},i}\big)}
             {v_{\mathrm{blk},i}}}
  \;\ge\; 1,
  \qquad
  g_{i}\equiv 1 \ \text{ where } v_{\mathrm{blk},i}=0 .
\end{equation}
Here $v_{\mathrm{blk}}=\diag\big(\sum_{V}C_{b}\big)$ is the diagonal of the
vertical block sum and $v^{\mathrm{cv}}_{\mathrm{uni}}$ is the diagonal of the
unified-throw covariance \emph{centered on the central value}. The bins where
the denominator vanishes are pinned to $g=1$ by declaration rather than
computed, and their count is carried in the product's receipt, because
``pinned'' and ``computed to one'' are different facts. The three band groups
$V$ (vertical), $R$ (residual) and $A$ (active lateral) are required to be
disjoint and exhaustive over the 45 bands. Implementation:
\texttt{nd-unfolding/z\_assembly.py}; the same algebra is asserted
independently of the producer that wrote the object.

\paragraph{Centering, explicitly.} The per-band blocks $C_{b}$ use the
MAT-conformant mean-centered $\pm1\sigma$ estimator of \S\ref{app:pair}. The
\emph{inflation operand} does not: variant \texttt{cv} takes
$v^{\mathrm{cv}}_{\mathrm{uni}}$, and the two centerings of the unified-throw
covariance differ by the outer product of the measured joint mean shift
$\mathbf{ms}$ (joint throw mean minus central value),
\begin{equation}
  \label{eq:zcentering}
  C^{\mathrm{cv}}_{\mathrm{uni}}
  = C^{\mathrm{mean}}_{\mathrm{uni}} + \mathbf{ms}\,\mathbf{ms}^{\top}
  \qquad\Longrightarrow\qquad
  v^{\mathrm{cv}}_{\mathrm{uni},i}-v^{\mathrm{mean}}_{\mathrm{uni},i}
  = \mathbf{ms}_{i}^{2},
\end{equation}
so the CV-centered variant retains the mean shift as variance where the
mean-centered variant discharges it. \textbf{That retention is the adopted
choice and is deliberate}: the mean-centered sibling \texttt{z-mean.npz} is
structurally identical --- same seven keys, same shapes, byte-identical central
values, support mask, pinned mask and row index --- and measures
$\sqrt{\Tr C}$ \SI{7.13}{\percent} \emph{lower}. The variant is therefore read
from the product's own metadata and checked, never inferred from a path:
\texttt{project\_cov\_nd.py -{}-expect-variant} is required with no default,
and \texttt{-{}-run-class publication} from a mean-centered source is refused.

\paragraph{The projection to $(\Eavail,W)$.} The quoted 42-cell covariance is
$C_{EW}=M\,C_{Z}\,M^{\top}$ with $M$ of shape
$\zprojCells\times\latBins$, built by
\texttt{project\_cov\_nd.build\_projection} over the kept axes
$(\Eavail,W)$. Each entry of $M$ is the product of the \emph{dropped} axes'
bin widths, so the destination is a differential density in the kept axes and
carries units \si{cm^2/GeV^2/nucleon}; converting a destination cell to an
integrated cross section requires multiplying by that cell's own
$\Delta\Eavail\,\Delta W$. Row $r$ of the result is destination dense-grid
index $\texttt{hRowIndex}[r]$ on the $7\times6$ $(\Eavail,W)$ grid in
\textbf{C order} --- measured, not assumed. Against the figure's central
values the C-order reshape agrees to a relative $1.7\times10^{-16}$ while the
F-order control disagrees at relative $1.0$, i.e.\ completely; the two orderings
are distinguished by the measurement rather than by convention. No source cell
is dropped ($\texttt{src\_cells\_dropped}=0$) and the result is exactly
symmetric ($\max|C-C^{\top}|=0$), a structural consequence of the
$M C M^{\top}$ form rather than a tolerance.

\paragraph{Conditioning, and why no rank is quoted.} Two properties of this
matrix are routinely misstated, including in an earlier version of this
subsection, so both are given with their operands.

\emph{It has no negative eigenvalue, and that fact is weaker than it sounds.}
The measured extremes are $\lambda_{\min}=+4.359104\times10^{-92}$ ---
\emph{positive} --- and $\lambda_{\max}=1.488215\times10^{-77}$, so
$\lambda_{\min}/\lambda_{\max}=2.9\times10^{-15}$ and the condition number is
$\approx3.4\times10^{14}$, at the edge of double precision. The commonly
transcribed label ``most-negative over max'' for that ratio is a misnomer:
there is no negative eigenvalue, and $n_{\mathrm{negative}}=0$. But the
smallest eigenvalues sit at the floating-point noise floor of a
$\latBins$-term weighted sum, so \textbf{their sign carries no information},
and $n_{\mathrm{negative}}=0$ is a statement about this computation --- not a
demonstration that the source's negative directions were cured, and not a
property the object need retain under a different mask, summation order or
precision.

\emph{The numerical rank is a property of a cutoff, not of the matrix.} The 42
eigenvalues decay smoothly over about fifteen orders of magnitude with no
plateau anywhere, so the retained count tracks whatever threshold is applied:
26 at a relative cutoff $10^{-6}$, 36 at the $10^{-12}$ hardcoded in
\texttt{project\_cov\_nd.py}, 41 at \texttt{numpy.linalg.matrix\_rank}'s own
default $n\epsilon=9.3\times10^{-15}$, and 42 at no cutoff at all.
\textbf{A bare ``rank 36'' should not be quoted, and it is not a number of
degrees of freedom.} What is true without a cutoff is qualitative: the
spectrum is continuous down to the noise floor, so any use requiring an
inverse requires a \emph{declared} regularization, which this analysis does
not supply silently. The cutoff scan is recorded in
\texttt{PLAN-20260918-scalar5d-publication-completion.md} \S16.1a, performed
on the diagnostic projection of the same source through the same $M$; its
$\lambda_{\min}$, $\lambda_{\min}/\lambda_{\max}$ and rank-at-$10^{-12}$ are
identical to the values the publication run recorded, and no separate scan on
the publication bytes exists.

\paragraph{Masks, and what the pairing is.} The support mask, pinned mask and
row-index arrays travel with the covariance and are byte-identical between the
\texttt{cv} and \texttt{mean} variants. The pairing between this covariance and
the central values shown in Fig.~\ref{fig:eavailWexcess} is established
\emph{by digest identity} --- the projected source's central-value array and the
figure's are the same bytes --- and not by numerical agreement; the numerical
agreement is a corroborating check.

\paragraph{What the adoption did and did not change.} The object's own metadata
is preserved unedited and still reads
$\texttt{scientific\_acceptance}=\text{\texttt{NON-PASSING}}$ and
$\texttt{adoptable}=\text{\texttt{false}}$. \textbf{An adoption does not turn
those fields into pass labels.} What was adopted is a decision to publish under
a digest-bound exception, recorded on 2026-09-20; the specification requires the
null bound to be fixed before production and production had already happened,
so the derivation step is unperformable for this object by anyone, which is why
this is an exception and not a repair. The exception attaches to these bytes and
to no other candidate.

\subsection{The four measurements that qualify the adopted covariance}
\label{sec:eavailw-adopted}
The adopted five-axis covariance is \texttt{z-cv.npz}
(\texttt{sha256}~\zcvDigest\ldots, \num{890500272} bytes, variant
\texttt{cv}), constructed as in \S\ref{sec:eavailw-construction}, and its
$(\Eavail,W)$ projection is \texttt{sha256}~\zprojDigest\ldots\ on
\zprojCells\ cells.

Four measurements travel with it. They are measurements, not caveats: each is a
number from a committed receipt, and each bounds what may be claimed.

\begin{enumerate}
\item \textbf{The estimator-seed sensitivity exceeds its declared limit.} A
  two-member campaign at estimator-seed offsets $\{0,1200\}$ graded the
  correlation-sensitive leg at $s_{\mathrm{proj}}=\sprojMeasured\,\%$ against a
  \sprojBound\,\% movement limit fixed before production. The two
  diagonal-only legs passed ($s_{\mathrm{agg}}=\saggMeasured\,\%$,
  $s_{\mathrm{med}}=\smedMeasured\,\%$): what the failing leg sees and they do
  not is off-diagonal structure, which is precisely what a projection uses.
\item \textbf{That sensitivity behaves like a property of the estimator rather
  than of the ensemble's resolution.} Measured with the throw ensemble held
  \emph{fixed} and only the seed changed, it does not fall with ensemble size
  over the range tested: mean $\seedEffectNforty\,\%$ over four disjoint
  40-throw subsets, $\seedEffectNeighty\,\%$ over the two 80-throw halves, and
  $\sprojMeasured\,\%$ on the full $N=160$ ensemble, giving a fitted exponent
  $\seedEffectExponent$ in $s\propto N^{-p}$. The same-seed resampling floor,
  measured at $N=40$ and $N=80$ only, falls steeply over that step ---
  $\floorNforty\,\%$ to $\floorNeighty\,\%$, a two-point exponent
  $\floorExponent$ --- and no floor measurement at $N=160$ exists. Statistical
  noise must fall with $N$; this did not fall between $N=40$ and $N=160$.
  \emph{Three limits on the reading, and they are the reason no
  infinite-ensemble statement is made here.} The $40$- and $80$-throw points are
  \emph{nested subsets} of the same 160 throws, not independent ensembles, so
  the flatness is a within-ensemble statement over a factor of four in $N$ and
  does not by itself determine the behaviour at much larger $N$; the
  $\pm\seedEffectScatter\,\%$ quoted with the mean is scatter across those throw
  subsets at a \emph{single} seed pair (offsets $\{0,1200\}$), so the width of
  the seed-pair distribution is unmeasured; and the statistic is a maximum over
  the declared offset set, so further pairs can only raise it. What is
  established is that the failure of the $\sprojBound\,\%$ limit is not the
  statistic's own resampling noise, and that no reduction was observed over
  $N=40$--$160$.
\item \textbf{Five bands lie outside what the variation probes.} Five of the
  covariance's 45 bands --- the muon-energy, muon-resolution and beam-angle
  blocks --- are produced at a fixed estimator seed that the offset does not
  reach, by construction of the endpoint convention. They carry
  $\pinnedBandsSqrtTr\,\%$ of $\sqrt{\mathrm{Tr}\,C}$ and contribute exactly zero
  movement. The variation therefore does not probe their seed sensitivity, and
  what the total would show were they also to vary is unmeasured \emph{in
  either direction}: $s_{\mathrm{proj}}$ is a maximum of a relative change over
  a set of functionals, not a sum of nonnegative component contributions, so
  releasing a component that was held fixed can move the total either way.
\item \textbf{The seed also moves the central values, by a small fraction of
  their own uncertainty.} On the 43 projection functionals of the grade (these \zprojCells\ destinations plus the all-ones total rate), the central value
  moves by at most $\cvMoveProjMax\,\%$, against relative uncertainties of order
  ten percent: \textbf{at most $\cvMoveOverSigmaMax\,\%$ of the quoted
  uncertainty, median $\cvMoveOverSigmaMed\,\%$}. The largest movement falls on
  the same destination cell that drives the covariance result. Per individual
  five-axis bin it reaches $\cvMoveBinOverSigmaMax\,\%$ of that bin's
  uncertainty; the projections do not inherit that, because summing over bins
  suppresses it.
\end{enumerate}

\noindent The predeclared magnitude criterion for this estimator-seed
sensitivity --- \emph{cause~3}'s $M(\mathrm{ii})$ in the internal
dependency register --- remains \emph{predeclared and not computed} for this
object: the 2026-09-20 campaign computed it for two \emph{other} products,
built for that purpose, and grading them does not regrade these bytes. The first, second and fourth measurements above were all made on those two products; only the third is a property of these bytes. The
exact record identifiers are in App.~\ref{app:provenance}.

The complementary cross-check against MINERvA's published
low-recoil measurement is presented in the $(\Eavail,q_{3})$ plane
(\S\ref{sec:4d}): on the maximal common grid the central values agree, while
no covariance-based consistency metric is reported.
